\section{Interval fields and a candidate full upper bound}
\label{S:interval-upper}

\myCmts{Full-target upper-bound candidate, pending independent review of
the interval coefficient argument and the source-limit transfer.
The accepted static results do not by themselves certify this section.}

Write
\[
  \nu=3-\alpha>1,\qquad
  \beta=4-\alpha-\alpha_0>2,\qquad p=2-\alpha_0>1.
\]
Thus $\beta-\nu=1-\alpha_0>0$.
In this section the spatial dimension is any admissible $d\in\{1,2,3\}$.
The spatial covariance is the nonnegative, even, positive-definite tempered
Radon measure $\Gamma$, with scaling
$\Gamma(rA)=r^{d-\alpha}\Gamma(A)$ and spectral convention
$\Gamma=\int e^{i\xi\cdot x}\mu(\ud\xi)$.
This includes $\Gamma=\delta_0$ when $d=\alpha=1$.
Set
\[
  R_\alpha=\frac{\alpha\pi\mu(B_1)}{2\sin(\pi\alpha/2)},\qquad
  A_\rho=\rho^2-\Delta,\qquad
  \|(f,g)\|_\rho^2=\|f\|_2^2+\|A_\rho^{-1/2}g\|_2^2.
\]
Complex Hilbert inner products are linear in the first slot.

\begin{lemma}[Dimension-independent Gaussian contractions]
  \label{L:interval-contractions}
  Uniformly in positive spatial smoothing and in finite Gaussian-mode
  and spatial Fourier-band cutoffs,
  \begin{equation}\label{E:interval-column}
    \sum_m\|A_\rho^{-1/2}M_{e_m}f\|_2^2
    \le R_\alpha\rho^{\alpha-2}\|f\|_2^2.
  \end{equation}
  Here $W_\epsilon(x)=\sum_m g_me_m(x)$ is an orthonormal-coordinate
  representation of $W(p_\epsilon(x-\cdot))$, whose covariance is
  $\Gamma_\epsilon=p_{2\epsilon}*\Gamma$.
  Consequently, on a Gaussian color that is initially in its vacuum,
  the phase insertion $(u,v)\mapsto(0,W_\epsilon u)$ has norm at most
  $\sqrt{R_\alpha}\rho^{(\alpha-2)/2}$ into its Gaussian $L^2$ phase.
  With Hermite degree of this color truncated at $N$, its sandwiched
  field has operator norm at most
  \begin{equation}\label{E:interval-fock}
    2\sqrt{N+1}\sqrt{R_\alpha}\rho^{\alpha/2-2}.
  \end{equation}
  Both statements remain valid tensoring with any other Hilbert factor.
\end{lemma}

\begin{proof}
  Fourier uniqueness gives $\mu(B_r)=r^\alpha\mu(B_1)$.
  The Gaussian comparison in Lemma~\ref{L:phase-contractions} uses no
  dimension-specific identity:
  \[
    \int e^{-s|\eta+\xi|^2}\mu(\ud\xi)
    =\int p_{2s}^{(d)}(x)e^{i\eta\cdot x}\Gamma(\ud x)
    \le\int e^{-s|\xi|^2}\mu(\ud\xi).
  \]
  Evenness gives a real integral, and physical nonnegativity gives
  the inequality. Integrating against $e^{-\rho^2s}\ud s$ and using
  radial spectral homogeneity yields
  \[
    \sup_\eta\int\frac{\mu(\ud\xi)}{\rho^2+|\eta+\xi|^2}
    \le R_\alpha\rho^{\alpha-2}.
  \]
  No exclusion of a physical origin atom is used.

  Gaussian isometry and unnormalized spatial Plancherel give
  \[
    \sum_m\|A_\rho^{-1/2}M_{e_m}f\|_2^2
    =\frac1{(2\pi)^d}\int|\widehat f(\eta)|^2
    \left(\int\frac{e^{-\epsilon|\xi|^2}\mu(\ud\xi)}
    {\rho^2+|\eta+\xi|^2}\right)\ud\eta .
  \]
  This proves~\eqref{E:interval-column}. A finite mode sum is a
  positive sub-sum. For a Fourier-band projection $S_b$, which commutes
  with $A_\rho$, use $S_bA_\rho^{-1}S_b\le A_\rho^{-1}$ to obtain
  the same inequality for $S_bM_{e_m}S_b$.
  The phase assertion on a new color follows because
  $\sum_mg_m(0,M_{e_m}u)$ has squared norm given by the left side
  of~\eqref{E:interval-column}, and its input is independent of this color.

  Substituting $A_\rho^{-1/2}f$ gives
  \[
    \sum_m(A_\rho^{-1/2}M_{e_m}A_\rho^{-1/2})^2
    \le R_\alpha\rho^{\alpha-4}I.
  \]
  These spatial sandwiches are selfadjoint, so their row and column
  variances agree. The normalized Hermite creation-row norm from
  degree $j$ to degree $j+1$ is $\sqrt{j+1}$, and annihilation is
  its adjoint, as proved in Lemma~\ref{L:fock-form}.
  Composition with the spatial column and addition of the two
  parts prove~\eqref{E:interval-fock}. The same positive operator
  identities hold with an identity on every other Hilbert factor.
  This proves Lemma~\ref{L:interval-contractions}.
\end{proof}

\begin{lemma}[Complex sums and changing phase weights]
  \label{L:interval-complex}
  The complex phase conclusion of Lemma~\ref{L:complex-phase} holds
  for $L=A_0-V$ whenever $V$ is bounded, not necessarily selfadjoint,
  \[
    \|A_\rho^{-1/2}VA_\rho^{-1/2}\|\le D\rho^{a-2},\qquad \rho>0,
    \quad 0<a<1,\quad D\le\delta_a\Lambda^{2-a}.
  \]
  Its bound is $100e^{2\Lambda h}$ in $X_\Lambda$.
  If a piecewise-constant potential is the sum of connected-interval
  fields, assign a positive weight $\omega_i$ to each field present
  and put $\Lambda(t)=\lambda+\sum_i\omega_i\mathbf1_{I_i}(t)$.
  The total loss from changing between these phase norms, starting
  and ending in $X_\lambda$, is at most
  \begin{equation}\label{E:interval-metric}
    \prod_i(1+\omega_i/\lambda).
  \end{equation}
\end{lemma}

\begin{proof}
  The adjoint has the identical sandwich bound, since $A_\rho$ is
  selfadjoint. Put $K_\kappa=A_\kappa^{-1/2}VA_\kappa^{-1/2}$ and
  $\kappa=(\Lambda^2+t)^{1/2}$. The resolvent difference is exactly
  \[
    (L+\Lambda^2+t)^{-1}-A_\kappa^{-1}
    =A_\kappa^{-1/2}(1-K_\kappa)^{-1}K_\kappa A_\kappa^{-1/2}.
  \]
  Its product on the left with $A_\Lambda^{1/2}$ has norm at most
  $D\kappa^{a-3}/(1-\delta_a)$. The inverse-root integral in
  Lemma~\ref{L:complex-phase} therefore gives the same $1/6$ bound.
  Apply it also to the adjoint and take inverses to obtain all four
  root comparisons, each with norm at most two.

  The real-part bound at $\rho=\Lambda$ is
  $\operatorname{Re}\langle(L+\Lambda^2)u,u\rangle
  \ge(1-\delta_a)\langle A_\Lambda u,u\rangle$.
  Separately, optimizing the form estimate at
  $\rho=\|A_0^{1/2}u\|/\|u\|$ bounds the imaginary part as in that
  lemma. Thus the same numerical-range parabola, scalar functions
  $C_h,F_h,H_h$, and four weighted phase blocks give $100e^{2\Lambda h}$.
  Selfadjointness of $V$ was not used; the adjoint estimate replaced it.

  For a common phase vector,
  \[
    \|z\|_{\rho_{\rm new}}
    \le\max\{1,\rho_{\rm old}/\rho_{\rm new}\}\|z\|_{\rho_{\rm old}}.
  \]
  At a jump, write $E$ for ending fields, $S$ for starting fields,
  and $R$ for those remaining. The ratio is at most
  \[
    \frac{\lambda+\sum_R\omega_i+\sum_E\omega_i}
    {\lambda+\sum_R\omega_i+\sum_S\omega_i}
    \le1+\sum_E\omega_i/\lambda
    \le\prod_E(1+\omega_i/\lambda).
  \]
  An initial increase of the weight has no loss. Count a final return
  to $X_\lambda$ as an ending jump at $T$.
  Each connected interval ends once. This proves
  \eqref{E:interval-metric} and Lemma~\ref{L:interval-complex}.
\end{proof}

\begin{lemma}[A positive inventory product]\label{L:inventory-product}
  For $k\ge1$, define
  \[
    v_i=i^{-\beta}+k^{-1}i^{-(\beta-1)}
    +k^{\alpha-2}i^{-p},\qquad
    b_{i,l}=\frac{i^{-\beta l}(1+li/k)^2}{(l!)^\nu},\quad l\ge2.
  \]
  There is a finite $C=C(\alpha,\alpha_0)$, independent of $M,k$, such that
  \begin{equation}\label{E:inventory-product}
    [z^k]\prod_{i=1}^M\left(1+v_i z+\sum_{l\ge2}b_{i,l}z^l\right)
    \le\frac{C^k}{(k!)^\nu}.
  \end{equation}
\end{lemma}

\begin{proof}
  Evaluate the positive product at $z=\theta k^\nu$, with $\theta=1/4$.
  Put $P_i=\theta k^\nu i^{-\beta}$ and
  $Q_i=\theta k i^{-p}$. The third singleton term is $Q_i$ because
  $\nu+\alpha-2=1$. The inequality $1+A+B\le(1+A)(1+B)$
  separates the factors $1+Q_i$.
  For $p>1$,
  \[
    \sum_{i\ge1}\log(1+A i^{-p})
    \le\int_0^\infty\log(1+A x^{-p})\ud x
    =A^{1/p}\int_0^\infty\log(1+x^{-p})\ud x.
  \]
  The last integral is finite. Hence their logarithmic contribution
  is at most $Ck^{1/p}\le Ck$.

  If $i\le k$, the remaining factor is at most
  $S_\nu(4P_i)=\sum_{l\ge0}(4P_i)^l/(l!)^\nu$:
  its singleton is at most $2P_i$, and
  $(1+li/k)^2\le(1+l)^2\le4^l$.
  Since $\nu>1$,
  $S_\nu(x)\le e^{\nu x^{1/\nu}}$.
  Therefore
  \[
    \sum_{i\le k}\log S_\nu(4P_i)
    \le\nu(4\theta)^{1/\nu}k\sum_{i\ge1}i^{-\beta/\nu}
    \le Ck,
  \]
  because $\beta/\nu>1$.

  If $i>k$, then $P_i\le\theta$. The remaining singleton tail sums to
  \[
    \theta\left(\frac1{\beta-1}+\frac1{\beta-2}\right)
    k^{\nu+1-\beta}=Ck^{\alpha_0}.
  \]
  Also, uniformly for $0\le P\le1/4$ and $y\ge0$,
  $\sum_{l\ge2}P^l(1+ly)^2/(l!)^\nu\le CP^2(1+y^2)$.
  Summing this bound gives
  \[
    Ck^{2\nu}\left(\frac{k^{1-2\beta}}{2\beta-1}
    +\frac{k^{-2}k^{3-2\beta}}{2\beta-3}\right)
    =Ck^{2\alpha_0-1}\le Ck.
  \]
  Use $\log(1+y)\le y$ for these fine factors.
  The entire product is consequently at most $e^{Ck}$.
  Positivity of its coefficients bounds the $k$th coefficient by
  $\theta^{-k}k^{-\nu k}e^{Ck}\le C^k/(k!)^\nu$.
  This proves Lemma~\ref{L:inventory-product}.
\end{proof}

Let $\mathcal F$ be a finite family of connected intervals $I_i$
meeting $[0,T]$, with lengths $h_i$ sorted decreasingly.
Let $W_i^\epsilon$ be independent copies of the smoothed unit field,
and suppose
\begin{equation}\label{E:interval-inventory}
  h_i\le L T/i,\qquad \sigma_i^2\le c h_i^{-\alpha_0}.
\end{equation}
Write $\mathcal U_{\mathcal F}^\epsilon(T;s)$ for phase evolution with
potential $s\sum_i\sigma_i\mathbf1_{I_i}(t)W_i^\epsilon(x)$.

\begin{proposition}[Uniform interval coefficients]\label{P:interval-coefficients}
  There is $C=C(\alpha,\alpha_0,L,c)$ such that, for real Schwartz phase $v$,
  every $k\ge1$ satisfies
  \begin{equation}\label{E:interval-coefficient-bound}
    \left|[s^{2k}]\E\|\mathcal U_{\mathcal F}^\epsilon(T;s)v\|_{1/T}^2\right|
    \le\frac{C^k R_\alpha^k T^{\beta k}}{(k!)^\nu}\|v\|_{1/T}^2,
  \end{equation}
  uniformly in $\epsilon>0$ and the finite family.
  In particular,
  \begin{equation}\label{E:interval-phase-bound}
    \E\|\mathcal U_{\mathcal F}^\epsilon(T;1)v\|_{1/T}^2
    \le e\exp\{C R_\alpha^{1/\nu}T^{\beta/\nu}\}\|v\|_{1/T}^2.
  \end{equation}
\end{proposition}

\begin{proof}
  Zero noise is trivial, and zero $\sigma_i$ can be omitted.
  First retain finitely many Gaussian coordinates and a bounded spatial
  Fourier band. For a fixed finite family and positive smoothing,
  the instantaneous field variance is bounded by some $v_{\max}<\infty$.
  On total Gaussian degree at most $n$, unweighted multiplication
  has norm at most $2\sqrt{n+1}\sqrt{v_{\max}}$.
  The $n$th phase Duhamel word is thus bounded by
  \[
    e^{\lambda T/2}
    \frac{(2\sqrt{n+1}\sqrt{v_{\max}}T/\lambda)^n}{n!}\|v\|_\lambda .
  \]
  This gives an entire $L^2$ vacuum series in the finitely many amplitudes.
  The bound need not be uniform in the inventory or smoothing.
  It justifies coefficient expansions, absolute energy double series,
  and cutoff identification below.

  Give every interval its own amplitude. The energy coefficient is a sum of
  $\E\langle U_{\boldsymbol r}v,U_{\boldsymbol s}v\rangle$, where
  $U_{\boldsymbol r}$ is the corresponding vector Duhamel coefficient.
  Gaussian reflection in each field makes this zero if $r_i+s_i$ is odd.
  Put $r_i+s_i=2k_i$ and $\sum_i k_i=k$. Inactive fields have $k_i=0$
  and may be deleted exactly. There are at most $k$ active fields and
  at most $\prod_i(2k_i+1)\le3^k$ splits of their degrees between copies.
  Use local Hermite cutoff $N_i=2k_i$.
  Each partial word in either copy has at most $2k_i$ appearances of
  that field, so this preserves the exact vacuum coefficient, including
  all annihilations and self-pairings.

  Set $\lambda=k/T$.
  A field with $k_i=1$ is either a mean field, with split $(2,0)$ or $(0,2)$,
  or a cross field, with split $(1,1)$.
  Fields with $k_i\ge2$ are called higher fields.
  Mean and higher fields will be estimated by Cauchy, while cross fields
  will be extracted linearly.
  For each Cauchy field choose
  \[
    \omega_i=k_i/h_i,\qquad
    d_i=2\sqrt{2k_i+1}\sqrt{R_\alpha},\qquad
    r_i^*=\frac{\delta_a\omega_i^{2-a}}{\sigma_i d_i},
    \qquad a=\alpha/2.
  \]
  Apply vector-valued Cauchy separately on the two phase copies.
  An unused amplitude is set equal to zero. In particular a mean field
  appears in only one copy.

  On either copy, hold the Cauchy amplitudes on these circles and set the
  cross amplitudes equal to zero. The resulting background is piecewise
  constant. With
  $\Lambda(t)=\lambda+\sum_{\text{Cauchy fields present}}\omega_i
  \mathbf1_{I_i}(t)$, its all-$\rho$ sandwich coefficient is at most
  \[
    \sum_i|\zeta_i|\sigma_i d_i
    \le\delta_a\sum_i\omega_i^{2-a}
    \le\delta_a\Lambda(t)^{2-a}.
  \]
  Lemmas~\ref{L:interval-contractions} and \ref{L:interval-complex}
  therefore bound each constant segment in $X_{\Lambda(t)}$ by
  $100e^{2\Lambda(t)\Delta t}$.
  There are at most $2k+1$ background segments. Splitting at the cross
  insertion times adds at most $k$ pieces. Moreover,
  \[
    \int_0^T\Lambda(t)\ud t
    \le\lambda T+\sum_i\omega_i h_i\le2k .
  \]
  All fixed segment constants and exponential factors on both copies are
  bounded by $C^k$. The metric-change factor on each copy is, by
  \eqref{E:interval-metric},
  $\prod_i(1+k_i/(\lambda h_i))$ over the Cauchy fields present there.
  Thus a mean field costs one such factor in energy, and a higher field
  costs at most its square.

  We spell out the simultaneous linear extraction of the cross fields.
  Its coefficient is the integral over
  $\prod_{i\in\mathrm{cross}}(I_i\cap[0,T])$ of the background-propagator
  word with its insertions sorted by their times.
  The permutation simplices partition this cube up to null diagonals;
  they do not introduce an additional factorial.
  Before its sole insertion, a cross field is in its Gaussian vacuum.
  The background and all other colors act as identity on this Gaussian
  factor. Its phase insertion is therefore the creation map of
  Lemma~\ref{L:interval-contractions}, at weight $\Lambda(t_i)$.
  Already-created other colors can be included in the remaining Hilbert
  factor. These embeddings compose with squared norm factors
  \[
    \sigma_i^2 R_\alpha\Lambda(t_i)^{\alpha-2}
    \le\sigma_i^2 R_\alpha\lambda^{\alpha-2}.
  \]
  Each background segment and each metric change occurs only once along
  a word; dividing a segment at an insertion adds a fixed constant,
  not a new power of its earlier costs.
  Minkowski over the cube gives the time factor $\prod_i h_i$ per copy.
  Cauchy--Schwarz between the two copies thus gives, for each cross field,
  the weight $R_\alpha\sigma_i^2h_i^2\lambda^{\alpha-2}$.
  This uses no annihilation bound on an already-occupied cross color:
  such a color is inserted only once on each copy.

  The resulting vectors are holomorphic in the remaining amplitudes.
  Cauchy's formula supplies total power $(r_i^*)^{-2k_i}$ for every
  mean or higher field. Since $d_i^2\le12k_iR_\alpha$ and $k_i!\le k_i^{k_i}$,
  \[
    (r_i^*)^{-2k_i}
    \le\frac{(C R_\alpha\sigma_i^2h_i^{4-\alpha})^{k_i}}
    {(k_i!)^\nu}.
  \]
  The final change from $X_\lambda$ to $X_{1/T}$ costs at most $k^2$
  in the energy, while the initial $X_\lambda$ norm is no larger.
  Absorb this polynomial and all split/segment constants into $C^k$.
  For a fixed variance multiindex, the absolute coefficient is bounded
  by the product of
  \[
    C R_\alpha\left[
    h_i^\beta(1+1/(\lambda h_i))
    +\lambda^{\alpha-2}h_i^{2-\alpha_0}\right],\quad k_i=1,
  \]
  and
  \[
    \frac{(C R_\alpha)^{k_i}h_i^{\beta k_i}}
    {(k_i!)^\nu}
    (1+k_i/(\lambda h_i))^2,\quad k_i\ge2 .
  \]
  The singleton expression includes both mean orientations and the
  cross case. Gaussian sign symmetry was used only to exclude odd
  total degrees, not to assume positivity of these phase coefficients.

  These weights are increasing in $h_i$. Substituting
  \eqref{E:interval-inventory} and $\lambda=k/T$, and extracting
  $C^kR_\alpha^kT^{\beta k}$, reduces their sum to the positive product
  in Lemma~\ref{L:inventory-product}. Factors depending on $L,c$ are
  absorbed per total variance order into $C^k$.
  This proves~\eqref{E:interval-coefficient-bound} at the cutoffs.

  At fixed positive smoothing and finite inventory the finite-degree
  multiplication bound using $v_{\max}$ is uniform in Gaussian-mode
  and spatial-band cutoffs. Gaussian-coordinate tails vanish on
  finite-coordinate vectors by dominated convergence; the annihilation
  tail is zero there, while the creation tail is bounded by the omitted
  pointwise variances. Density and closedness of actual multiplication
  identify the limits. Fourier projections converge strongly and commute
  with free evolution. Dominated time integration gives every fixed word,
  and the displayed $n!$ bound gives uniform convergence on compact
  amplitude sets. Thus the coefficients above are those of the actual
  finite-interval model. No zero-smoothing limit is taken here.

  The zero coefficient is at most $e\|v\|_{1/T}^2$.
  Since $\nu>1$, $\sum_{k\ge0}x^k/(k!)^\nu\le e^{\nu x^{1/\nu}}$.
  Absolute energy summation proves~\eqref{E:interval-phase-bound},
  after enlarging $C$. This proves
  Proposition~\ref{P:interval-coefficients}.
\end{proof}

\begin{theorem}[Candidate full upper scale]\label{T:full-upper}
  Under Hypothesis~\ref{H:homog}, for every admissible $d=1,2,3$,
  there is a finite $C=C(\alpha,\alpha_0)$ such that
  \begin{equation}\label{E:full-upper-bound}
    \E u^2(T,x)
    \le e(T+2)^d
    \exp\{C R_\alpha^{1/(3-\alpha)}
    T^{(4-\alpha-\alpha_0)/(3-\alpha)}\},\qquad T\ge1.
  \end{equation}
  In particular, this would establish the unchanged Goal~\ref{G:upper}.
\end{theorem}

\begin{proof}
  Zero noise is trivial. Choose $H_0$ with $6T\le H_0<12T$.
  At lengths $h_j=H_0 2^{-j}$, use the three grids with endpoints
  $h_j(n+(-1)^j t)$, $n\in\mathbb Z$, $t=0,1/3,2/3$.
  They are nested: the difference between the parent and child shifts
  in child units is an integer, since $3t\in\mathbb Z$.
  For $s\ne r$ in $[0,T]$, choose a scale
  $3|s-r|<h_j\le6|s-r|$.
  Boundaries of the three grids at this scale have mutual spacing
  $h_j/3>|s-r|$. Hence at least one grid interval contains both points,
  including their endpoints after choosing a boundary-free grid.
  It follows that
  \begin{equation}\label{E:dyadic-time-majorant}
    |s-r|^{-\alpha_0}
    \le6^{\alpha_0}\sum_{\text{three grids},\,|I|\le H_0}
    |I|^{-\alpha_0}\mathbf1_I(s)\mathbf1_I(r)=:K_\infty(s,r).
  \end{equation}
  The sum is at most $C_{\alpha_0}|s-r|^{-\alpha_0}$ off the diagonal:
  only lengths at least $|s-r|$ contribute, with at most one interval
  per grid and length. Endpoint values do not affect the time integrals.

  Keep intervals meeting $[0,T]$ through depth $J$, and use independent
  static smoothed fields with $\sigma_I^2=6^{\alpha_0}|I|^{-\alpha_0}$.
  Their time covariance is the corresponding partial sum $K_J$.
  The decreasing interval lengths obey $h_i\le L T/i$ with an absolute
  $L$. Indeed, at length $h_j$ there are at most $3(T/h_j+2)$ intervals,
  and summing coarser levels bounds their cumulative count by
  $L T/h_j$, using $H_0/T<12$.
  Proposition~\ref{P:interval-coefficients} applies uniformly in $J,\epsilon$.

  To pass from phase data to constant initial data, choose
  $f_T\in C_c^\infty$ with $0\le f_T\le1$, equal to one on $B_{T+1}$
  and supported in $B_{T+2}$, and set the velocity to zero.
  For $|y|+s<T+1$, $G_s*f_T(y)=s$, so $C_sf_T(y)=1$.
  The support of $G_t$ is contained in $B_t$.
  Induction on Duhamel order therefore identifies the position generated
  from $(f_T,0)$ at time $T$ on $B_1$ with the constant-data model
  $u_{J,\epsilon}$.
  This identity can be taken in Gaussian $L^2$ term by term.
  At finite $J,\epsilon$, pointwise convergence for constant data follows
  from the bound
  $\sqrt{(2n-1)!!}\,v_{\max}^{n/2}T^{2n}/(2n)!$
  on the $n$th position term.
  The model is spatially stationary. Consequently,
  \[
    |B_1|\E|u_{J,\epsilon}(T,0)|^2
    \le\E\|u_{f_T,J,\epsilon}(T)\|_2^2
    \le e\,|B_{T+2}|\exp\{C R_\alpha^{1/\nu}T^{\beta/\nu}\}.
  \]

  We finally compare pairing integrals, not arbitrary-phase Gaussian
  covariance orders. From the source
  \cite{chen:author2026:time-dependency}*{(2.23)--(2.24), p.~13},
  its actual common noise regularization has covariance
  $\gamma_\delta^0(s-r)\Gamma_\epsilon(x-y)$, where
  \[
    0\le\gamma_\delta^0(t)
    =\Gamma(\alpha_0)^{-1}\int_0^{\delta^{-1}}
    e^{-\lambda|t|}\lambda^{\alpha_0-1}\ud\lambda
    \le |t|^{-\alpha_0}.
  \]
  Here $\Gamma(\alpha_0)$ denotes the Euler Gamma function.
  At fixed noise smoothing, the source's Green-kernel limit is the
  classical measure-kernel Duhamel integral: bounded continuous spatial
  covariances and Gaussian moment bounds permit weak Green convergence
  before the noise limit, as in Section~\ref{S:G0}.

  Every two-chain Wick integrand for constant data is nonnegative in
  $d=1,2,3$, including self-pairings and unequal chain lengths.
  For a fixed finite partial-order sum of the source regularizations,
  \eqref{E:dyadic-time-majorant} bounds its second moment by the same
  finite-order pairing sum with $K_\infty\Gamma_\epsilon$.
  Nonnegative Tonelli makes this the increasing limit of the sums with
  $K_J\Gamma_\epsilon$. Each is bounded by the full constant-data moment
  just estimated, whose additional pairing terms are nonnegative.
  Thus the bound is uniform in the original noise parameters and in
  the partial-order cutoff. This uses only entrywise covariance
  domination inside positive Wick integrals; no coupling or Gaussian
  convex-order assertion is made.

  The spectral condition (1.5) required for the source construction
  follows by radial spectral homogeneity and $\alpha+\alpha_0<2$.
  From its Theorem~4.1, (4.2), p.~22, each fixed-order noise-regularized
  term converges in $L^2$ to $X_n$. Take that limit at a fixed partial
  order. Then let the partial order tend to infinity using
  Proposition~\ref{P:mixed} and Lemma~\ref{L:crude-corrected}, which
  establish the original $L^2$ series independently of this bound.
  The volume ratio is $|B_{T+2}|/|B_1|=(T+2)^d$.
  This proves~\eqref{E:full-upper-bound} and
  Theorem~\ref{T:full-upper}.
\end{proof}

\begin{remark}[Claim and calibration boundary]\label{R:full-upper-boundary}
  This section is a proof candidate until its complete source-bound review.
  It does not adopt a sharp constant, a first-moment normalization repair,
  a novelty verdict or formal verification.
  Its exponent tends to $(4-\alpha)/(3-\alpha)$ as $\alpha_0\downarrow0$.
  Cauchy--Schwarz and the unconditional Jensen inequality remain compatible
  with the nonsharp upper constant; any numeric $L_*$ calibration retains
  the source-normalization qualification stated earlier.
  The Skorohod exponent is compared, not used to identify constants:
  both Gaussian creation and annihilation were retained here.
  Goal~\ref{G:sharp} and the private Lean companion remain separate.
\end{remark}
